% TOP_PROBLEM: {"id": 30004807, "title": "Weak Bianchi Identities Across Timelike Singularities", "category_id": 16, "category": {"id": 16, "name": "physics", "display_name": "Mathematical Physics", "description": "Problems at the intersection of mathematics and physics.", "slug": "physics", "order_index": 16, "created_at": "2026-07-31T15:26:25.671Z"}, "status": "open", "statusQualification": "The cited work proves the weak twice-contracted second Bianchi identity for a restricted static spherically symmetric class. Its extension to all admissible timelike singularities is the remaining conjecture. Integrability and the treatment of the singular locus are part of the question; the classical pointwise identity alone does not justify this weak formula."}
% FIRST_POSTED: 2026-10-11
% ENTRY_KIND: statement_only
% SUBMISSION: {"kind": "statement_only", "category": "top_problems", "problem_id": "30004807", "proposal_type": "established", "source": "Mathematical Aspects of General Relativity (2021), Oberwolfach Reports, DOI 10.4171/owr/2021/40, pp. 2157--2267. Original report and mathematical context. https://doi.org/10.4171/owr/2021/40", "status_evidence": "Source status retained from the supplied canonical registry. Dated literature assessment (2026-08-22): The cited work proves the weak twice-contracted second Bianchi identity only for a restricted static spherically symmetric class; the extension to all admissible timelike singularities remains a conjecture. Primary report inspected for formulation; no comprehensive current-status literature review is claimed.", "status_checked_on": "2026-10-11", "duplicate_check": "Canonical numeric ID and public code checked against full registry; absent from baseline selected website registry. Source scope compared with adjacent source questions and existing related catalogue entries. No new identity allocated.", "attribution": "Original authors and contributors credited in Mathematical Aspects of General Relativity (2022), Oberwolfach Reports, DOI 10.4171/owr/2021/40, pp. 2157--2267; UnsolvedMath curation credited separately. Definitions and English exposition prepared with Codex assistance; no independent human mathematical verification claimed."}
% SOURCE_STATEMENT: "Does the weak contracted second Bianchi identity $\\int_M(R_{\\mu\\nu}-\\tfrac12R g_{\\mu\\nu})\\nabla^\\mu\\psi^\\nu\\,dV_g=0$ hold for every spacetime whose timelike singularity restricts on spacelike slices to a regular zero-area singularity with negative ZAS mass in the sense of Bray and Jauregui, and for every compactly supported vector field $\\psi$?"
% Preparation and final self-audit: CONTRIBUTING.md prompts C and D; source/audit artifacts retained in the local physics_statements_2026_10_11 batch.
% Imported from: physics_statements.tex; UnsolvedMath ID 30004807
% !TeX program = lualatex
\documentclass[11pt]{article}
\usepackage[margin=25mm]{geometry}
\usepackage{fontspec}
\setmainfont{Latin Modern Roman}
\usepackage{amsmath,amssymb,mathtools}
\usepackage{unicode-math}
\setmathfont{Latin Modern Math}
\usepackage{xurl,hyperref}
\hypersetup{colorlinks=true,urlcolor=blue}
\setcounter{secnumdepth}{0}
\setlength{\parindent}{0pt}
\setlength{\parskip}{5pt}
\setlength{\emergencystretch}{3em}
\newcommand{\sref}[2]{\hyperlink{source:#1:#2}{[S#2]}}
\newcommand{\cataloguescope}[1]{\paragraph{Recorded target.}#1}
\newcommand{\R}{\mathbb R}
\newcommand{\C}{\mathbb C}
\newcommand{\N}{\mathbb N}
\newcommand{\Z}{\mathbb Z}
\newcommand{\Q}{\mathbb Q}
\newcommand{\F}{\mathbb F}
\newcommand{\eps}{\varepsilon}
\begin{document}
% UnsolvedMath ID 30004807; source catalogue code OWR-8415343-010
\setcounter{section}{30004806}
\section{Weak Bianchi Identities Across Timelike Singularities}\label{30004807}
{\small UnsolvedMath ID 30004807 · Mathematical Physics}
\subsection{Definitions and mathematical statement}

\paragraph{Shared notation.}$\mathbb N=\{1,2,\ldots\}$; the standard number fields have their usual meanings. All additional conventions are stated locally.


The Einstein tensor is $G_{\mu\nu}=R_{\mu\nu}-\tfrac12R g_{\mu\nu}$, where $R_{\mu\nu}$ and $R$ are Ricci and scalar curvature. The derivative $\nabla$ and volume measure $dV_g$ are those of $g$. A weak identity means that the integral against every compactly supported smooth test vector field is well-defined and vanishes. ZAS denotes a zero-area singularity, with regularity and mass defined as in the report's Bray--Jauregui convention.

\paragraph{Catalogue wording retained for traceability.}
Does the weak contracted second Bianchi identity $\int_M(R_{\mu\nu}-\tfrac12R g_{\mu\nu})\nabla^\mu\psi^\nu\,dV_g=0$ hold for every spacetime whose timelike singularity restricts on spacelike slices to a regular zero-area singularity with negative ZAS mass in the sense of Bray and Jauregui, and for every compactly supported vector field $\psi$?

\paragraph{Scope and source qualification.}The cited work proves the weak twice-contracted second Bianchi identity for a restricted static spherically symmetric class. Its extension to all admissible timelike singularities is the remaining conjecture. Integrability and the treatment of the singular locus are part of the question; the classical pointwise identity alone does not justify this weak formula.
\subsection{Short English statement}
Decide whether the distributional conservation identity survives the stated class of timelike negative-mass zero-area singularities.
\subsection{Sources}
\begin{itemize}
\item[\hypertarget{source:30004807:1}{S1}] ULAM.AI and the UnsolvedMath Contributors, UnsolvedMath, record 30004807 (OWR-8415343-010). Catalogue curation under CC BY 4.0; definitions and English exposition added locally.
\url{https://huggingface.co/datasets/ulamai/UnsolvedMath}.
\item[\hypertarget{source:30004807:2}{S2}] Mathematical Aspects of General Relativity (2021), Oberwolfach Reports, DOI 10.4171/owr/2021/40, pp. 2157--2267. Original report and mathematical context.
\url{https://doi.org/10.4171/owr/2021/40}.
\end{itemize}
\label{end:30004807}
\end{document}
